\documentclass[11pt,a4paper]{article}

\def\nyear {2025}
\def\nterm {Winter}
\def\nlecturer {}
\def\ncourse {Introduction to Numerical Analysis}
\def\ncoursehead {Numerical Analysis}

\makeatletter

% packages
\usepackage{amssymb,amsfonts,amsmath,calc,tikz,pgfplots,geometry,mathtools}
\usepackage{color}   % May be necessary if you want to color links
\usepackage[svgnames]{xcolor}
\usepackage[hidelinks]{hyperref}
\usepackage{forest}
\usepackage{commath} % ew
\usepackage{esvect} % \vv makes vectors
\usepackage{centernot} % for the comparison
\usepackage{esdiff}
\usepackage{esint}
\usepackage{amsthm}
\usepackage{fancyhdr}
\usepackage{bm}
\usepackage{witharrows}
\usepackage{bookmark}
\usepackage{tikz-cd}
\usepackage{quiver}
\usepackage{bbm}
\usepackage{textcomp}
\usepackage{float}
\usepackage{gensymb}
\usepackage{cleveref}
\usepackage{environ}
\usepackage{comment}
\usepackage{cancel}
\usepackage{xparse}

% Inevitable cs
\usepackage{listings}
\usepackage[]{algorithm2e}

% tikz libraries
\usetikzlibrary{positioning}
\usetikzlibrary{matrix}
\usetikzlibrary{arrows}
\usetikzlibrary{bending}
\usetikzlibrary{arrows.meta}
\usetikzlibrary{decorations.markings}
\usetikzlibrary{decorations.pathreplacing}
\usetikzlibrary{decorations.markings}
\usetikzlibrary{patterns}
\usetikzlibrary{backgrounds}

% Page style setup
\pagestyle{fancy}
\fancyhfoffset{0pt}
\renewcommand{\headwidth}{\textwidth}
\geometry{margin=1in}
\pgfplotsset{compat=1.18}
\setlength{\headheight}{14.6pt}
\addtolength{\topmargin}{-1.6pt}
\hypersetup{
    colorlinks=false,
    linktoc=section,
    linkcolor=black,
}

%% maketitle setup
\ifx \nauthor\undefined
  \def\nauthor{Yeheli Fomberg}
\else
\fi

\ifx \nemail\undefined
  \def\nemail{\href{mailto:yp@campus.technion.ac.il}
  {\texttt{<yp@campus.technion.ac.il>}}}
\else
\fi

\ifx \ncoursehead \undefined
  \def\ncoursehead{\ncourse}
\fi

\lhead{\emph{\nouppercase{\leftmark}}}
\ifx \nextra \undefined
  \rhead{
    \ifnum\thepage=1
    \else
      \ncoursehead
    \fi}
\else
  \rhead{
    \ifnum\thepage=1
    \else
      \ncoursehead \ (\nextra)
    \fi}
\fi

\def\notes{notes}
\def\problems{problems}

\ifx \ntype \undefined
  \def\ntype{notes}
\else
\fi

\ifx \ntype \notes
  \let\@real@maketitle\maketitle
  \renewcommand{\maketitle}{\@real@maketitle\begin{center}
  \begin{minipage}[c]{0.9\textwidth}\centering\footnotesize
  These notes are not endorsed by the lecturers.
  I have revised them outside lectures to incorporate supplementary
  explanations, clarifications, and material for fun.
  While I have strived for accuracy, any errors or misinterpretations 
  are most likely mine.
  \end{minipage}\end{center}}
\else
  \ifx \ntype \problems
    \let\@real@maketitle\maketitle
    \renewcommand{\maketitle}{\@real@maketitle\begin{center}
    \begin{minipage}[c]{0.9\textwidth}\centering\footnotesize
    These problems are mostly based on the problem sets I was given
    when taking the course.
    
    While I have strived for accuracy and completeness,
    some of the problems would not have been here without the help of
    friends.
    \end{minipage}\end{center}}
  \else
  \fi
\fi


% theorem environments
\theoremstyle{definition}
\newtheorem{definition}{Definition}[section]
\newtheorem{remark}{Remark}[section]
\newtheorem{example}{Example}[section]
\newtheorem{exercise}{Exercise}[section]
\newtheorem{paradox}{Paradox}[section]
\newtheorem{entry}{Entry}[section]
\newtheorem*{solution}{Solution}
\newtheorem*{question}{Question}
\newtheorem*{answer}{Answer}
\newtheorem*{problem}{Problem}
\theoremstyle{plain}
\newtheorem{theorem}{Theorem}[section]
\newtheorem{proposition}[theorem]{Proposition}
\newtheorem{lemma}[theorem]{Lemma}
\newtheorem{corollary}[theorem]{Corollary}
\newenvironment{proofof}[1]{%
    \renewcommand{\proofname}{Proof of #1}%
    \begin{proof}
}{%
    \end{proof}
}
% tikz customization
\pgfarrowsdeclarecombine{twolatex'}{twolatex'}{latex'}{latex'}{latex'}{latex'}
\tikzset{->/.style = {decoration={markings,
                                  mark=at position 0.999
                                  with {\arrow[scale=2]{latex'}}},
                      postaction={decorate}}}
\tikzset{<-/.style = {decoration={markings,
                                  mark=at position 0.001 with {\arrowreversed[scale=2]{latex'}}},
                      postaction={decorate}}}
\tikzset{-->/.style = {decoration={markings,
                                  mark=at position 0.999
                                  with {\arrow[scale=2]{latex'[sep=0pt -2.5]}}},
                      postaction={decorate}}}
\tikzset{<--/.style = {decoration={markings,
                                  mark=at position 0.001
                                  with {\arrowreversed[scale=2]{latex'[sep=0pt -2.5]}}},
                      postaction={decorate}}}
\tikzset{<->/.style = {decoration={markings,
                                  mark=at position 0.001 with {\arrowreversed[scale=2]{latex'}},
                                  mark=at position 0.999 with {\arrow[scale=2]{latex'}},},
                      postaction={decorate}}}
\tikzset{->-/.style = {decoration={markings,
    mark=at position 0.50+0.225/#1 with {\arrow[scale=2]{latex'}}},
                      postaction={decorate}}}
\tikzset{-<-/.style = {decoration={markings,
                                   mark=at position 0.50-0.225/#1 with {\arrowreversed[scale=2]{latex'}}},
                       postaction={decorate}}}
\tikzset{->>/.style = {decoration={markings,
                                  mark=at position 1 with {\arrow[scale=2]{latex'}}},
                      postaction={decorate}}}
\tikzset{<<-/.style = {decoration={markings,
                                  mark=at position 0 with {\arrowreversed[scale=2]{twolatex'}}},
                      postaction={decorate}}}
\tikzset{<<->>/.style = {decoration={markings,
                                   mark=at position 0 with {\arrowreversed[scale=2]{twolatex'}},
                                   mark=at position 1 with {\arrow[scale=2]{twolatex'}}},
                       postaction={decorate}}}
\tikzset{->>-/.style = {decoration={markings,
                                   mark=at position 0.50+0.450/#1 with {\arrow[scale=2]{twolatex'}}},
                       postaction={decorate}}}
\tikzset{-<<-/.style = {decoration={markings,
                                   mark=at position 0.50-0.450/#1 with {\arrowreversed[scale=2]{twolatex'}}},
                       postaction={decorate}}}

\pgfarrowsdeclare{biggertip}{biggertip}{%
  \setlength{\arrowsize}{1pt}
  \addtolength{\arrowsize}{.1\pgflinewidth}
  \pgfarrowsrightextend{0}
  \pgfarrowsleftextend{-5\arrowsize}
}{%
  \setlength{\arrowsize}{1pt}
  \addtolength{\arrowsize}{.1\pgflinewidth}
  \pgfpathmoveto{\pgfpoint{-5\arrowsize}{4\arrowsize}}
  \pgfpathlineto{\pgfpointorigin}
  \pgfpathlineto{\pgfpoint{-5\arrowsize}{-4\arrowsize}}
  \pgfusepathqstroke
}
\tikzset{
	EdgeStyle/.style = {>=biggertip}
}

\tikzset{circ/.style = {fill, circle, inner sep = 0, minimum size = 3}}
\tikzset{scirc/.style = {fill, circle, inner sep = 0, minimum size = 1.5}}
\tikzset{mstate/.style={circle, draw, black, text=black, minimum width=0.7cm}}

\tikzset{eqpic/.style={baseline={([yshift=-.5ex]current bounding box.center)}}}

\definecolor{mblue}{rgb}{0.2, 0.3, 0.8}
\definecolor{morange}{rgb}{1, 0.5, 0}
\definecolor{mgreen}{rgb}{0, 0.4, 0.2}
\definecolor{mred}{rgb}{0.5, 0, 0}

% topology
\DeclareMathOperator{\dfr}{dfr} % deformation retract
\newcommand{\Cells}{\text{Cells}}
\newcommand{\RP}{\mathbb{RP}}
\newcommand{\CP}{\mathbb{CP}}

% order theory
\newcommand{\join}{\vee}
\newcommand{\meet}{\wedge}

% algebraic geometry
\DeclareMathOperator{\Rad}{Rad} % Radical of modules
\DeclareMathOperator{\Spec}{Spec}
% \renewcommand{\div}{\mathrm{div}} % divisors
\DeclareMathOperator{\Mer}{Mer} % Meromorphic functions
\DeclareMathOperator{\Exc}{Exc} % Exceptional divisor
\DeclareMathOperator{\Td}{Td} % Todd classes

% algebra

\DeclareMathOperator{\trdeg}{tr.deg.}
\DeclareMathOperator{\odd}{odd}
\DeclareMathOperator{\even}{even}
\DeclareMathOperator{\lcm}{lcm}
\DeclareMathOperator{\ord}{ord}
\DeclareMathOperator{\codim}{codim}
\DeclareMathOperator{\Ann}{Ann}
\DeclareMathOperator{\Ext}{Ext}
\DeclareMathOperator{\Sym}{Sym}
\DeclareMathOperator{\Out}{Out}
\DeclareMathOperator{\Aut}{Aut}
\DeclareMathOperator{\End}{End}
\DeclareMathOperator{\Inn}{Inn}
\DeclareMathOperator{\Mat}{Mat}
\DeclareMathOperator{\Fix}{Fix}
\DeclareMathOperator{\std}{std}
\DeclareMathOperator{\sgn}{sgn}
\DeclareMathOperator{\adj}{adj}
\DeclareMathOperator{\id}{id}
\DeclareMathOperator{\op}{op}
\DeclareMathOperator{\tr}{tr}
\DeclareMathOperator{\diag}{diag}
\DeclareMathOperator{\rank}{rank}
\DeclareMathOperator{\rk}{rk}
\DeclareMathOperator{\coker}{coker}
\DeclareMathOperator{\Mult}{Mult} % multiplier algebra
\DeclareMathOperator{\Frac}{Frac} % fraction field
\DeclareMathOperator{\Rat}{Rat}
\DeclareMathOperator{\Gal}{Gal}
\newcommand{\ioplus}{\overset{\text{int}}{\oplus}} % interior direct sum
\newcommand{\GG}{\mathcal G} % Galois correspondence
\newcommand{\FF}{\mathcal F} % Galois correspondence
\newcommand{\cupp}{\smile} % cup product
\newcommand{\capp}{\frown} % cap product
\newcommand{\actson}{\curvearrowright}
\newcommand{\bra}{\left\langle}
\newcommand{\ket}{\right\rangle}
\newcommand{\idealin}{\triangleleft\,}
\newcommand{\normalin}{\triangleleft\,}
\newcommand{\ip}[2]{\left\langle #1, #2 \right\rangle}
\newcommand{\bigslant}[2]
{{\raisebox{.2em}{$#1$}\left/\raisebox{-.2em}{$#2$}\right.}}
\newcommand{\properideal}{%
  \mathrel{\ooalign{$\lneq$\cr\raise.22ex\hbox{$\lhd$}\cr}}}

% categories
\newcommand{\cat}[1]{\mathbf{#1}}
\newcommand{\Ch}{\cat{Ch}} % Chain complexes   
\newcommand{\Set}{\cat{Set}}
\newcommand{\Grp}{\cat{Grp}}
\newcommand{\Ab}{\cat{Ab}}
\newcommand{\Ring}{\cat{Ring}}
\newcommand{\CRing}{\cat{CRing}}
\newcommand{\Top}{\cat{Top}}
\newcommand{\Vect}{\cat{Vect}}
\newcommand{\cmod}{\cat{mod}}

%% matrix groups
\newcommand{\GL}{\mathrm{GL}}
\newcommand{\Or}{\mathrm{O}}
\newcommand{\PGL}{\mathrm{PGL}}
\newcommand{\PSL}{\mathrm{PSL}}
\newcommand{\PSO}{\mathrm{PSO}}
\newcommand{\PSU}{\mathrm{PSU}}
\newcommand{\SL}{\mathrm{SL}}
\newcommand{\msl}{\mathrm{sl}}
\newcommand{\SO}{\mathrm{SO}}
\newcommand{\Spin}{\mathrm{Spin}}
\newcommand{\Sp}{\mathrm{Sp}}
\newcommand{\SU}{\mathrm{SU}}
\newcommand{\U}{\mathrm{U}}
\let\char\relax
\newcommand{\char}{\text{char}}

% analysis
\renewcommand{\Re}{\mathrm{Re}}
\renewcommand{\Im}{\mathrm{Im}}
\NewDocumentCommand{\dx}{o}{
  \IfNoValueTF{#1}
    {\dif x}                  % if no optional argument
    {\dif^{#1} \!x}         % if optional argument
}
\NewDocumentCommand{\dy}{o}{
  \IfNoValueTF{#1}
    {\dif y}                  % if no optional argument
    {\dif^{#1} \!y}         % if optional argument
}
\newcommand{\dt}{\dif t}
\newcommand{\du}{\dif u}
\newcommand{\dv}{\dif v}
\newcommand{\dz}{\dif z}
\newcommand{\ds}{\dif s}
\newcommand{\dw}{\dif w}
\newcommand{\dr}{\dif r}
\newcommand{\dm}{\dif m}
\newcommand{\df}{\dif f}
\newcommand{\dV}{\dif V}
\newcommand{\dS}{\dif S}
\newcommand{\dA}{\dif \mathrm{A}}
\newcommand{\dmu}{\dif \mu}
\newcommand{\dnu}{\dif \nu}
\newcommand{\dtheta}{\dif \theta}
\newcommand{\domega}{\dif \omega}
\newcommand{\dsigma}{\dif \sigma}
\newcommand{\dtau}{\dif \tau}
\newcommand{\dvarphi}{\dif \varphi}
\renewcommand{\dif}{\mathrm{d}}
\renewcommand{\Dif}{\mathrm{D}}
\renewcommand{\triangle}{\Delta}
\newcommand{\ptriangle}{\partial \triangle}
\DeclareMathOperator{\im}{im}
\DeclareMathOperator{\cis}{cis}
\DeclareMathOperator{\rot}{rot}
\DeclareMathOperator{\vspan}{span}
\DeclareMathOperator{\grad}{grad}
\DeclareMathOperator{\curl}{curl}
\DeclareMathOperator{\esssup}{esssup}
\newcommand{\hodge}{{\mathnormal{\star}}}
\DeclareMathOperator{\range}{range}
\DeclareMathOperator{\Hom}{Hom}
\DeclareMathOperator{\Int}{Int}
\DeclareMathOperator{\Conv}{Conv} % convex hull
\newcommand{\ind}{\mathbbm{1}}
\DeclareMathOperator{\Res}{Res}
\DeclareMathOperator{\graph}{graph}
\DeclareMathOperator{\diam}{diam}
\DeclareMathOperator{\supp}{supp}
\DeclareMathOperator{\Vol}{Vol} % Volume
\DeclareMathOperator{\Area}{Area} % Area
\DeclareMathOperator{\area}{area}
\DeclareMathOperator{\Ind}{Ind} % Index
\newcommand{\length}{\mathrm{length}}

% logic
\DeclareMathOperator{\MOD}{MOD}
\DeclareMathOperator{\Theory}{Theory}
\newcommand{\notimplies}{%
  \mathrel{{\ooalign{\hidewidth$\not\phantom{=}$\hidewidth\cr$\implies$}}}}

% nice
\newcommand{\half}{\frac{1}{2}}
\newcommand{\pair}{\del}
\newcommand{\taking}[1]{\xrightarrow{#1}}
\newcommand{\otaking}[1]{\xleftarrow{#1}}
\newcommand{\inv}{^{-1}}
\newcommand{\ot}{\leftarrow}
\newcommand{\ninfty}{-\infty}
\newcommand{\pinfty}{+\infty}
\newcommand{\floor}[1]{\left\lfloor #1 \right\rfloor}
\newcommand{\ceil}[1]{\left\lceil #1 \right\rceil}
\newcommand{\viff}{\Big\Updownarrow} % vertical iff
\newcommand{\bmid}{\,\middle\vert\,} % big mid

% probability
\newcommand{\Prob}{\mathbf{P}}
\renewcommand{\vec}[1]{\boldsymbol{\mathbf{#1}}}
\DeclareMathOperator{\Bin}{Bin}
\DeclareMathOperator{\Geo}{Geo}
\DeclareMathOperator{\Poi}{Poi}
\DeclareMathOperator{\Exp}{Exp}
\DeclareMathOperator{\Var}{Var} % Variance
\DeclareMathOperator{\Cov}{Cov}

% special letters
\newcommand{\N}{\mathbb{N}}
\newcommand{\Z}{\mathbb{Z}}
\newcommand{\Q}{\mathbb{Q}}
\newcommand{\R}{\mathbb{R}}
\newcommand{\C}{\mathbb{C}}
\newcommand{\F}{\mathbb{F}}
\newcommand{\E}{\mathbf{E}}
\newcommand{\D}{\mathbb{D}}
\renewcommand{\H}{\mathbb{H}}
\newcommand{\ps}{\mathcal{P}}
\newcommand{\M}{\mathcal{M}}
\renewcommand{\L}{\mathcal{L}}
\renewcommand{\O}{\mathcal{O}}
\renewcommand{\U}{\mathcal{U}}
\newcommand{\Omicron}{O}
\newcommand{\powerset}{\mathcal{P}}

% text
\newcommand{\st}{\text{ s.t. }}
\newcommand{\tand}{\quad \text{and} \quad}
\newcommand{\tor}{\quad \text{or} \quad}
\newcommand{\stand}{\text{ and }}
\newcommand{\stor}{\text{ or }}
\renewcommand{\tt}[1]{\textnormal{\textbf{(#1).}}} %tt=theorem title GET RID OF

% title format
\title{\textbf{\ncourse}}

\ifx \ntype \notes
  \author{Notes taken by \nauthor\, \nemail \\ Based on lectures by \nlecturer}
  \date{\nterm\ \nyear}
\else
  \ifx \ntype \problems
    \author{These problems were solved by \nauthor \\ \nemail}
  \else
  \fi
\fi

\makeatother


\DeclareMathOperator{\fl}{fl}
\DeclareMathOperator{\NaN}{NaN}

\begin{document}
\maketitle

% Insert cool image here

\newpage
\tableofcontents
\newpage

\section{Introduction}
This course addresses what of all the math we have learned so far can
we compute with the computer

For example, given the task of computing the determinant of a matrix $A$,
we want to find the most efficient algorithm.
The most efficient algorithm is the one that uses the least amount of
operations.

One way of calculating $\det A$ is using permutations:
\[
  \det A = \sum_{\gamma \in S_n}
  \sgn(\gamma) a_{1 \sigma(1)} \cdot a_{2 \sigma(2)} \cdots a_{n \sigma(n)}.
\]
Each permutation costs us $n$ operations of multiplication, and to calculate
the sum we perform $n! - 1$ summation operations.
Since there are $n!$ permutations in $S_n$, the total amount of operations
of this algorithm is $n \cdot n! + n! - 1$.

An alternative algorithm would be to find the eigendecomposition
$A = U^{T} D U$ where $U$ is an orthogonal matrix, and then calculate
the product of the elements on the main diagonal of $D$ (which are the
eigenvalues of $A$).
Using this algorithm we can compute $\det A$ in $C n^3 + n - 1$ operations
for some costant $C$.
Although this algorithm is much faster than the previous ones, there are
more efficient algorithms still.

\begin{definition}[Big $\O$ notation]
  Let $(x_n)_{n \geq 1}$ and $(y_n)_{n \geq 1}$ be sequences of real numbers.
  We say that $x_n = \O(y_n)$ if there exist constants $C$ and $N$ such that
  for all $n > N$
  \[
    |x_n| \le C |y_n|
  \]
\end{definition}

Here are some examples
\begin{align*}
  \frac{n + 1}{n^2} &= \O\del{\frac 1n} \\
  Cn^3 + n - 1 &= \O(n^3) \\
  n \cdot n! + n! - 1 &= \O(n \cdot n!)
\end{align*}

Using big $\O$ notation we can say that direct determinant calculations
requires $\O(n \dots n!)$ calculations, while using eigendecomposition
to calculate it takes $\O(n^3)$ operations.

However, the big $\O$ notation is still not very satisfying because
we still have
\[
  C n^3 + n - 1 = \O(n \cdot n!).
\]
To fix this problem we introduce the $\Theta$ notation.

\begin{definition}[big $\Theta$ notation]
  Let $(x_n)_{n \geq 1}$ and $(y_n)_{n \geq 1}$ be sequences of real numbers.
  We say that $x_n = \Theta(y_n)$ if there exist constants $0 < c < C$ and $N$ 
  such that for all $n > N$
  \[
    c |y_n| \le |x_n| \le C |y_n|.
  \]
\end{definition}

Here are some examples
\begin{align*}
  a_k n^k + \cdots + a_1 n + a_0 &= \Theta(n^k) \\
  Cn^3 + n - 1 &= \Theta(n^3) \\
  n \cdot n! + n! - 1 &= \Theta(n \cdot n!)
\end{align*}

\begin{example}[Horner's algorithm]
  Another important method before we move to the next section,
  is Horner's method.
  It allows computing a polynomial $p(x) = \sum_{k=0}^{n} a_k x^n$ in
  $\Theta(n)$ operations instead of $\Theta(n^2)$ operations.
  It states that
  \[
    \begin{aligned}
      &a_{0} + a_{1}x + a_{2}x^{2} + a_{3}x^{3} + \cdots + a_{n}x^{n} \\ = {}
      &a_{0} + x\del{a_{1} + x\del{a_{2} + x\del{a_{3} + \cdots + 
      x(a_{n-1} + x a_{n}) \cdots}}}.
    \end{aligned}
  \]
\end{example}

\section{Digital Number Representation}
We usually use base-$10$ expansion to represent number.
For example, the number $x = 123.45$ means
\[
  x = 1 \times 10^2 + 2 \times 10^1 + 3 \times 10^0 + 4 \times 10^{-1} +
    5 \times 10^{-2}.
\]
In this case, we say that $x$ has a finite base-$10$ expansion.
All non-negative real numbers have a (possibly infinite) base-$10$ expansion.

\begin{theorem}
  \label{thm:bases}
  Let $x$ be a real number in $[0,1)$, and $b \geq 2$ an integer.
  Then there exists $L \in \Z$ and a sequence $(c_k)_{k=1}^{\infty}$ where
  each $c_k$ is in $\set{0,1,\dots,b-1}$, so that
  \[
    x = \sum_{k=1}^{\infty} c_k b^{-k}.
  \]
\end{theorem}

The motivation for this theorem is to prove for example, that computers,
who work in base $2$, can represent the same numbers we can represent
in the way that is more convenient to us like base $10$.

\begin{lemma}
  \label{lem:bases}
  If $b \geq 2$ is an integer, $L \in \Z$ and $x \in [0,b^{L+1})$.
  Then there exists $c_L \in \set{0,1,\dots,b-1}$ such that
  $x - c_L b^{L} \in [0,b^L)$.
\end{lemma}
\begin{proof}
  We notice that
  \[
    \bigcup_{j=0}^{b-1} \intco{j b^L, (j + 1) b^L} =
    [0, b^{L+1}).
  \]
  So there exists $c_L \in \set{0,1,\dots,b-1}$ so that 
  $x \in [c_L b^L, (c_L + 1) b^L)$.
  Then
  \[
    r_L = x - c_L b^L \in [0,b^L).
  \]
\end{proof}

We can now go back to prove \Cref{thm:bases}

\begin{proof}
  By assumption $x \in [0,b^0)$.
  By \Cref{lem:bases} with $L = -1$ we now define:
  \[
    r_1 := x - c_1 b^{-1} \in [0,b^{-1}).
  \]
  And continue indefinitely
  \[
    r_S := x - \sum_{k=1}^{S} c_k b^{-k} \in [0,b^{-k}).
  \]
  This implies that $r_S \taking{S \to \infty} 0$ which means that
  \[
    x = \sum_{k=1}^{\infty} c_k b^{-k}.
  \]
\end{proof}

We can now conclude from the theorem that for $x > 0$, we can find $L$
large enough so that $x < b^L$ and so $x b^{-L} < 1$.
We can then write
\[
  x = b^L \sum_{k=1}^{\infty} c_k b^{-k}.
\]

It is also possible in base-$b$ for $b \geq 2$ to represent any
non-negative number as:
\[
  x = (-1)^s \times [1.f]_b \times 2^m,
\]
for some $s \in \set{0,1}$, $m \geq 0$ and $f$ an infinite sequence
of digits.

\begin{remark}
  Obsviouly we can also represent $0$ because $0 \in \set{0,1,\dots,b}$.
\end{remark}

\subsection{Representing numbers in the computer}

\begin{definition}[Integers]
  In a computer, we represent integers using $32$ bits.
  We use one bit to represent the sign of the integer,
  and the other as constants to the powers of two.
  In this way we can represent every integer 
  $n \in [- (2^{-31} - 1), 2^{31} - 1]$ as such:
  \[
    n = (-1)^{c_31} \sum_{j=0}^{30}.
  \]
  Where $(c_0,c_1,\dots,c_{31})$ represent the bits.
  If we add two numbers and get a result larger than $2^31 - 1$,
  the computer will give us the result $Inf$.
\end{definition}

\begin{definition}[Fixed point representation]
  One (uncommon) way of representing rational numbers with $32$ bits is
  called fixed point representation.
  In this method we allot the last bit for the sign,
  and we write the numbers as such
  \[
    x =
    (-1)^{c_{31}} (c_{30},\dots,c_{15}.c_{14},\dots,c_1,c_0)_{2} =
    (-1)^{c_{31}}
    \del{c_{30} 2^{15} + c_{29} 2^{14} +
    \cdots + c_{15} + \cdots + c_0 2^{-15}}.
  \]
\end{definition}

\begin{definition}[Floating point]
  The common way to represent rational numbers in the computer is by
  using the floating point format.
  In this method we write rational numbers as
  \[
    x = (-1)^s \times [1.f]_2 \times 2^m.
  \]
  We use a single bit to encode the sign, eight bits to encode $m$ (also
  known as the exponent) and the remaining $32$ bits to encode $f$ (also known
  as the normalized mantissa). We call $[1.f]_2$ the mantissa.
  This method is known as single precision.

  Precise details about the implementation of special number like
  $\pm 0$, $\pm 0$ or $\NaN$ are the in solutions.

  In double precision (which is the standard in Matlab) we use $64$ bits
  to store each number. A single bit to encode the sign, $11$ bits to encode 
  the exponent and the remaining $52$ bits to encode the normalized mantissa.
  In general we only focus on single precision.

  The largest floating point number in single precision is
  \[
    2^{127} \le
    x = (1.1 1 \cdots 1) \times 2^{127} \le
    2^{128}.
  \]

\end{definition}

\begin{definition}[Machine number]
  Any number that can be represented in the computer exactly is called
  a machine number.
  In most cases we refer to numbers that can be represented using the floating
  point method.
\end{definition}

Denote $\floor x$ or $\fl(x)$ the closest machine number to $x$.
\begin{definition}[Absolute error]
  The absolute error is defined as the value of $\abs{\fl(x) - x}$ and
  we can see that in the case of a floating point that
  \[
    \abs{\fl(x) - x} \le 2^m 2^{-24}
  \]
\end{definition}

\begin{definition}[Relative error]
  Let $f(x)$ be some approximation of $x \neq 0$, the relative error
  is defined by
  \[
    \frac{|f(x) - x|}{|x|} < \epsilon
  \]
  if and only if
  \[
    f(x)= x(1 + \delta)
  \]
  for some $\delta \in (-\epsilon, \epsilon)$.
  We can get this equlity by setting $\delta = \frac{f(x) - x}{x}$.
\end{definition}

\begin{remark}
  The relative erroor in the case of a floating point approximation satisfies
  \[
    \frac{\abs{\fl(x) - x}}{|x|} \le
    \frac{2^m 2^{-24}}{[1.f]_2 \times 2^m} \le
    2^{-24}.
  \]
\end{remark}
This means that the relative error of rounding is pretty low.
In double precision because we have more bits the relative error
is bounded by $2^{-52} \sim 2 \times 10^{-16}$.
This number is called $eps$ in Matlab.

\begin{definition}[Unit roundoff error]
  The unit roundoff error in a machine, is defined to be the minimal upper bound
  of the relative error.
  In single precision, the unit roundoff error is $2^{-24}$.
\end{definition}

\begin{definition}[Machine precision, machine epsilon]
  The machine precision (or machine epsilon) is defined to be the distance
  between the number $1$ and the immediate following machine number.
  In single precision, the machine epsilon is $2^{-23}$.
  That is because
  \[
    (1.0 0 \cdots 0 1)_2 = 1 + 2^{-23}.
  \]
  The closest number from below however is
  \[
    (1.1 1 \cdots 1 1)_2 \times 2^{-1} = 1 - 2^{-24}.
  \]
  In double precision the machine epsilon drops to $2^{-52}$.
\end{definition}

\begin{proposition}
  Let $x_1,x_2 > 0$ and assume there exists some $f \colon \R \to \R$ such that
  there exist $\delta_1$, $\delta_2$, $\delta_3$ such that
  \[
    f(x_1) = (1 + \delta_1) x_1 \stand
    f(x_2) = (1 + \delta_2) x_2 \stand
    f(f(x_1) \cdot f(x_2)) = (1 + \delta_3)[f(x_1) \cdot f(x_2)]
  \]
  where $|\delta_1|, |\delta_2|, |\delta_3| < \epsilon$.
  Then
  \[
    f(f(x_1) \cdot f(x_2)) =
    (1 + \delta) x_1 x_2
  \]
  For some $|\delta| < 3 \epsilon + 3 \epsilon^2 + \epsilon^3$.
\end{proposition}
\begin{proof}
  \[
    f(f(x_1) \cdot f(x_2)) =
    (1 + \delta_3)[f(x_1) \cdot f(x_2)] =
    (1 + \delta_3)(1 + \delta_2)(1 + \delta_1) x_1 x_2
  \]
  and
  \[
    (1 + \delta_3)(1 + \delta_2)(1 + \delta_1) =
    1 + \delta_1 + \delta_2 + \delta_3 + 
    \delta_1 \delta_2 + \delta_1 \delta_3 + \delta_2 \delta_3 + 
    \delta_1 \delta_2 \delta_3.
  \]
  Denote
  \[
    \delta =
    \delta_1 + \delta_2 + \delta_3 + 
    \delta_1 \delta_2 + \delta_1 \delta_3 + \delta_2 \delta_3 + 
    \delta_1 \delta_2 \delta_3.
  \]
  we get
  \[
    |\delta| \le 3 \epsilon + 3 \epsilon^2 + \epsilon^3
  \]
  which completes the proof.
\end{proof}

The above theorem shows that when taking the product of two numbers,
the error is bigger than epsilon, but not by much.
The situation is similar for addition and division,
but not for subtraction.

Consider
\[
  x_1 = [1.f0]_2 \stand x_2 = [1.f6]
\]
we have that $\floor{x_1} - \floor{x_2} = 0$ which implies that
relative error is
\[
  \frac{\abs{\floor{\floor{x_1} - \floor{x_2}} - |x_1 - x_2|}}{|x_1 - x_2|} = 1.
\]
This occurs when the distance between $x_1$ and $x_2$ is very small compared
to $x_1$.

\begin{proposition}
  Assume that $x_2 > x_1 > 0$ and $\frac{x_2}{x_2 - x_1} < a$ for some $a$.
  Let $f \colon \R \to \R$ be a function satisfiying
  \[
    f(x_1) = (1 + \delta_1) x_1 \stand
    f(x_2) = (1 + \delta_2) x_2 \stand
    f(f(x_1) - f(x_2)) = (1 + \delta_3)[f(x_1) - f(x_2)]
  \]
  where $|\delta_1|, |\delta_2|, |\delta_3| < \epsilon$.
  Then
  \[
    f(f(x_1) - f(x_2)) = (1 + \delta)(x_2 - x_1)
  \]
  where $|\delta| < 2 \epsilon + 2 a \epsilon + \epsilon^2 + 2 a \epsilon^2$.
\end{proposition}
\begin{proof}
  The proof is simply by calculation
\end{proof}

Because of the propositions we have seen so far, some ways to compute
certain expressions are more accurate than others.
For example instead of computing
\[
  x - \sin(x) \tor
  \sqrt{x^2 + 1} - 1
\]
to get a more accurate calculation we can calculate the equivalent
expressions
\[
  - \frac{x^3}{3!} + \frac{x^5}{5!} - \frac{x^7}{7!} \cdots
  \tand
  \frac{x^2}{\sqrt{x^2 + 1} + 1}.
\]

As mentioned earlier, in single precision we can not represent numbers bigger
than $2^{128}$ or smaller than $2^{-128}$.
Instead the computer will return us the values $\infty$ and $\ninfty$ 
respectively. This is called \emph{overflow}.
We also can not represent numbers in the range $(-2^{-128}, 2^{-128})$.
Instead the computer will return us the values $0$ or $-0$ if the number
was negative. This is called \emph{underflow}.

\begin{definition}[Softmax function]
  For a fixed $\alpha \geq 0$, the softmax function 
  $s^{\alpha} \colon \R^d \to \R^d$ is defined as
  \[
    s^{\alpha}_i := \frac{e^{\alpha x_i}}{\sum_{j=1}^{d} e^{\alpha} x_j}.
  \]
\end{definition}
\begin{remark}
  For all $x \in \R^d$ and $\alpha \geq 0$ the entries of $s^{\alpha}(x)$
  are all positive and sum to one.
\end{remark}
\begin{remark}
  If $k$ is some index such that $x_k > x_i$ for all $i \neq k$, then
  \[
    \lim_{\alpha \to \infty} s^{\alpha}(x) = e_k
  \]
  where $e_k$ is the unit vector with $1$ in the $k$th index and $0$ otherwise.
  Thus, for large $\alpha$ softmax gives a strong idication of where the
  maximum of the vector is located.
\end{remark}
\begin{remark}
  If $\alpha = 0$ then $s^0(0) = (1/d, 1/d, \dots, 1/d)$ for all $x$.
\end{remark}

This function is intresting because although the entries in the result vectors
are all between $0$ and $1$, the exponentiation can lead to underflow and
overflow which will greatly affect its computation.

To avoid overflow when computing softmax we can notice that for any scalar
$M$ and vector $1_d = (1, 1, \dots, 1)$ we have that 
$s^{\alpha}(x) = s^{\alpha}(x - M 1_d)$.
This allows use to choose $M = \max\set{x_i}_i$ and compute
$s^{\alpha}(x - M 1_d) = s^{\alpha}(x)$ without overflow.

Softmax is useful in machine learning in for \textit{classification tasks}.
For example, suppose you have $d$ different kanji characteres registered in
some database by order of their JIS code.
We want to build a function $h$ that assigns each vector $X$ representing a
picture of a single handwritten charachter a probability vector
\footnote{A vector with nonnegative entries that sum up to one.}
$(h_1(X),\dots,h_d(X))$ so that $h_i(X)$ represents the probability
of the picture to describe the $i$th character.
The standard way of doing this is by defining
\[
  h(X) := s^{\alpha} A_2 \sigma A_1 (X),
\]
where $A_1$, $A_2$ are matrices, $\sigma$ is a nonlinear transformation,
and $s^{\alpha}$ is the softmax function.

This is a simple neural network, where the paramter $\alpha$ controls how
''sharp`` the probability vector $h(X)$ will be.

The real challenge is how to define $h$ well,
but this is beyond the scope of this lecture.
Instead we can notice that underflow doesn't really affect softmax,
and while overflow does lead to complete nonsense,
we can make sure it never happens.
More on that in the solutions section.

\subsection{Alternatives for symbolic computation*}

The method we have used so far to calculate nubmbers is called numercial
computation.
Because of the way we represent numbers in (floating points) we always
get an error bounded by $2^{-24}$.
There are some alternatives to numerical calculation.

One alternative for numerical calculations is \emph{symbolic calculations}.
This method is used by Mathematica, Maple, Wolfram Alpha 
and can also be used in Matlab.
It has two main behaviours
\begin{itemize}
  \item When multiplying numbers, the numbers of digits increase to avoid
    truncation.
  \item Solutions of equations can be given symbolically like $\sqrt{2}$,
    $\pi$ etc.
\end{itemize}
Using this method we gain more accuracy, but lose speed of computation.

Another alternative is called \emph{interval arithmetics}.
Since our calculations are not exact we can represent numbers using intevals.
For example, let $x = [1.f]_2$ such that $f = f_1 \circ f_2$.
where $f_1$ is the first 23 digits of $f$ and $f_2$ is the rest of the digits.
In this case, we use the interval $I_x = [[1.f_1]_2,[1.f_1 1]_2]$ to represent
$x$ because $x \in I_X$.

Suppose we have two numbers $x = [a_x, b_x]$ and $y = [a_y, b_y]$, when
calculating the sum $x + y$ we get $x + y = [a_x + a_y, b_x + b_y]$.
In this way we can be certain that
$x + y \in I_{x+y} = [a_x + a_y, b_x + b_y]$.
We define multiplication in a similar, though more complicated way.
When tying to define multiplication for intercal arithmetics, one could
use the following proposition.

\begin{proposition}
  The minimum and maximum of the function $f(x, y) = xy$ over
  the rectangle $[a_x, b_x] \times [a_y, b_y]$ are obtained at the corners. 
\end{proposition}

\newpage

\section{Solving non-linear equations}
From now on, we assume we can store numbers without inaccuracies, and
perfrom arithmetics without errors.

\subsection{Bisection method}
Suppose we wanted to find a solution to the equation
\[
  \sin(x) = e^x.
\]
Then from analysis courses we know to approximate the solution by changing
the equaition to standard form
\[
  f(x) = e^x - \sin(x)
\]
and then find an interval like $[a, b] = [-\frac{3 \pi}{2}, 1]$ for which
$f(a) < 0 < f(b)$.
Because then we know that there exists a solution $x \in [a,b]$.

The bisection method allows us to find roots of continuous functions in
intervals $[a,b]$ such that $\mathrm{sign} f(a) \neq \mathrm{sign} f(b)$.
It works by repeating the following process
\begin{enumerate}
  \item[(-1)] Set $a_0 = a$, $b_0 = b$ and $c_0 = \frac{a_0 + b_0}{2}$.
  \item[(n)] Assume $a_n$, $b_n$, $c_n$ are defined. If $f(c_n) = 0$
    we have found a root. Otherwise, from the pigeonhole principle there
    exists $x \in \set{a_n,b_n}$ such that $f(c_n)$ and $f(x)$ form an interval
    that contains $0$. We call this interval $[a_{n + 1}, b_{n + 1}]$.
    We set $c_{n + 1} = \frac{a_{n + 1} + b_{n + 1}}{2}$. Treating $n$
    as a variable for formality, we increment it by one, and continue to
    the next step (step $(n)$).
\end{enumerate}
Setting $|a - b| = L$ we get that $|a_n - b_n| = L 2^{-n}$.
By Cantor's intersection theorem, and the squeeze theorem, it is clear
that $c_n \taking{n \to \infty} x$ where $x$ is a root of $f(x)$.

\begin{remark}
  Using the bisection method, we can calculate the number of steps we need to
  take to obtain $x'$ where $x' \in [x - \epsilon, x + \epsilon]$ for any
  $\epsilon > 0$.
\end{remark}

Here is pseudocode for the bisection method, $f$ represents a continuous
function, $a$, $b$ represent the limits of the interval $[a,b]$ where
the root is.
Tolerance is the word used for the accuracy we need, represented by $\epsilon$.
We also need to have a maximum amount of iterations denoted $N_{\text{max}}$.
\begin{algorithm}[H]
\SetAlgoLined
\KwIn{$f$, $a$, $b$, $\epsilon$, $N_{\text{max}}$}
\KwOut{Approximate root $c$ or failure to converge}

\If{$f(a) \cdot f(b) \geq 0$}{
    Print "Error: No root in the interval $[a, b]$"\;
    \Return
}

\For{$k \gets 1$ \textbf{to} $N_{\text{max}}$}{
    $c \gets \frac{a + b}{2}$\;
    \If{$f(c) = 0$ \textbf{or} $\frac{b - a}{2} \leq \epsilon$}{
        \Return $c$\;
    }
    \If{$f(a) \cdot f(c) < 0$}{
        $b \gets c$\;
    }
    \Else{
        $a \gets c$\;
    }
}

Print "Failed to converge after $N_{\text{max}}$ iterations"\;
\end{algorithm}

\subsection{Newton's algorithm}
Newton's algorithm (also known as the Newton--Raphson algorithm) is an algorithm
to compute the root of a function $f \colon \R \to \R$
so that $f \in C^2(\R)$ and $f$ has a root $x$.

The idea behind the algorithm is to have a first guess of the root $x_0$,
and then set $x_{n+1}$ to be the root of the Taylor expansion up to linear
order of $f$ at the point $x_n$.

\begin{remark}
  This algorithm does not always converge to the root,
  % example
  but if we have a ``good'' approximation for the root as $x_0$ it will always
  converge to the root very quickly.
  We will see what ``good'' exactly means later.
\end{remark}

\begin{example}
  Consider the function $f(x) = e^x - 6$.
  The solution is $x = \log 6 \approx 1.7918$.
  We start by a close guess $x_0 = 2$.
  The iterations are given by
  \[
    x_{n + 1} =
    x_n - \frac{f(x_n)}{f'(x_n)} =
    x_n - \frac{e^{x_n} - 6}{e^{x_n}}.
  \]
  Applying this process gives
  \[
    x_1 = 1.8120 \stand
    x_2 = 1.7920 \stand
    x_3 = 1.7918 \stand
    x_4 = 1.7918
  \]
  with errors
  \[
    e_1 = 2 \times 10^{-2} \stand
    e_2 = 2 \times 10^{-4} \stand
    e_3 = 2 \times 10^{-8} \stand
    e_4 = 4 \times 10^{-16}.
  \]
  which shows that the process is very useful.
\end{example}

\begin{lemma}
  \label{lem:newton-local}
  Let $f \in C^2(\R)$ be a function with a root $r \in \R$.
  Let $x_{n + 1}$ be the solution of $L(x \mid x_n) = 0$
  (the root of the linear Taylor approximation of $f(x)$ at $x_n$).
  Then there exists $\xi \in (r,x_n)$ such that
  \[
    f'(x_n)(x_{n+1} - r) = \half f''(\xi)(x_n - r)^2
  \]
\end{lemma}
\begin{proof}
  By definition we have $L(x_{n + 1} \mid x_n) = 0$, and Taylor's expansion
  around $x_n$ gives us
  \[
    0 = f(r) = L(r \mid x_n) + \half f''(\xi)(x_n - r)^2
  \]
  so
  \[
    f'(x_n)(x_{n+1} - r) = L(x_{n + 1} \mid x_n) - L(r \mid x_n) =
    L(r \mid x_n) = \half f''(\xi)(x_n - r)^2
  \]
  which completes the proof.
\end{proof}

\begin{proposition}
  Let $f \in C^2(\R)$ be a function with a simple root $r \in \R$.
  Then there exists $\delta > 0$ and $C > 0$ such that for any Newton
  iterations starting from $x_0 \in [r - \delta, r + \delta]$, the
  sequence $|x_n - r|$ converges monotonely to zero and
  \[
    |x_{n  + 1} - r| \le C(x_n - r)^2
  \]
\end{proposition}
\begin{remark}
  A simple root means that $f'(r) \neq 0$.
\end{remark}
\begin{proof}
  Let $x_n$ be the sequence Newton--Raphson iterates starting from $x_0$.
  Since $f'(r) \neq 0$, we can choose some $\delta_0 > 0$ small enough so
  that $f'(x) \neq 0$ on $U = [r - \delta_0, r + \delta_0]$.
  Set
  \[
    m = \min_{x \in U} |f'(x)| \stand M = \max_{x \in U} |f''(x)|
  \]
  then using \Cref{lem:newton-local} we have that whenever 
  
  to be added
\end{proof}

\begin{remark}
  The convergence rate of Newton's method near a minimun is qaudratic.
  This means that
  \[
    |x_{n + 1} - r| \le C (x_n - r)^2.
  \]
  In comparison, the convergence rate of the bisection method was linear
  \[
    |x_{n + 1} - r| \le \half |x_n - r|.
  \]
\end{remark}

\begin{definition}[Convergence rate]
  Let $x_n$ be a sequence converging to $L$.
  We say that $x_n$ has a convergence rate $\alpha > 0$ if there exists
  $C > 0$ and $N > 0$ such that for all $n < N$ we have
  \[
    |x_{n + 1} - L| \le C |x_n - L|^{\alpha}
  \]
\end{definition}

\begin{theorem}
  Assume $f \in C^2(\R)$ is strictly increasing, strictly convex, 
  and has a root.
  Then the root is unique, and Newton iterations will converge to it from
  any starting point.
\end{theorem}
\begin{proof}
  It is clear that the root is unique.
  From strictly convex, and strictly increasing we have that
  $f''(x) > 0$ and $f'(x) > 0$ for all $x$.
  From \Cref{lem:newton-local} we see that $x_n - r > 0$ for all $n \geq 1$.
  Since $f$ is increasing it follows that $f(x_n) > f(r) = 0$.
  Since
  \[
    x_{n + 1} = x_n - \frac{f(x_n)}{f'(x_n)}
  \]
  we see that $r < x_{n + 1} < x_n$, so $x_n$ is monotonically decreasing and
  bounded below, thus it converges to some limit $L$.
  Since $f \in C^2(\R)$ we have that
  \[
    L = 
    \lim_{x \to \infty} x_{n + 1} =
    \lim_{n \to \infty} x_n - \frac{f(x_n)}{f'(x_n)} =
    L - \frac{f(L)}{f'(L)}
  \]
  which means that $f(L) = 0$ and thus $L = r$ which completes the proof.
\end{proof}

\subsection{Generalization to higher dimensions}
Let $f \colon \R^d \to \R^d$ be a function such that $f \in C^2{\R^d}$
and $f$ has a root $r \in \R^d$.
Given a point $x_n \in \R^d$, we can define $x_{n + 1}$ as the solution
to the linear approximation of $f$ around $x_n$,
\[
  L(x \mid x_n) = f(x_n) + D_f(x_n)(x - x_n).
\]
The solution for the equation $L(x \mid x_n) = 0$ is
\[
  x_{n + 1} = x_n- D_f^{-1}(x_N) f(x_n).
\]
Given that $D_f$ is non-singular, we get a quadratic convergence rate.

\begin{remark}
  While the bisection method allowed linear convergence under very minimal
  convergence, Newton's method allows
  (at least \footnote{See solutions for cubic convergence of Newton's method
  in the solutions section}) quadratic convergence,
  and it can also be generalized to higher dimensions.
\end{remark}

\subsection{Newton's method with the contractive mapping theorem}
The iterations in Newton's method are given by
\[
  x_{n+1} = x_n - \frac{f(x_n)}{f'(x_n)}.
\]
Define a function $F$ as follows
\[
  F(x) = x - \frac{f(x)}{f'(x)}.
\]
The functions is well defined for points where the derivative of $f$ doesn't
vanish, and also satisfies
\[
  x_{n+1} = F(x_n) \tand F(r) = r \iff f(r) = 0.
\]
Algorithms that compute sequences like $x_{n+1} = F(x_n)$ aim to find a fixed
point of $F$ at the limit.
These algorithms are called fixed point interations.
One of the most helpful thoerems when dealing with these algorithms,
it the contractive mapping theorem.

\begin{theorem}[Contractive mapping theorem]
  Let $(X,d_X)$ be a completes metric space and
  let $F \colon X \to X$ be a contraction with contracion factor $\lambda$.
  Then $F$ has a unique fixed point $r \in X$.
  Moreover, every sequence $\set{x_n}_{n=1}^{\infty} \subseteq X$ of the
  form
  \[
    x_{n+1} = F(x_n)
  \]
  converges to $r$ at least linearly.
  Note that when working with an interval, the image of $[a,b]$ must
  be contained in $[a,b]$.
\end{theorem}
\begin{proof}
  To be added.
\end{proof}

% Analysis of algortithms

\begin{example}[Gradient descent]
  Another fixed point iteration method is the gradient descent method.
  Assume that $f \colon \R \to \R$ is an element of $C^2(\R)$,
  and let $r$ be a strict miniumum ($f'(r) = 0$ and $f''(r) > 0$).
  The iterations in the gradient descent method are defined as such
  \[
    F(x_n) = x_n - \alpha f'(x_n)
  \]
  where $\alpha > 0$ is called the \textit{step-size}.
\end{example}
\begin{remark}
  Each step we ``descend'' in the direction of the derivative.
\end{remark}
\begin{remark}
  This method also works in higher dimensions by changing the iterations
  to
  \[
    F(x_n) = x_n - \alpha \nabla f(x_n).
  \]
\end{remark}

\begin{example}
  Assume we want to minimize the function $f(x) = x^2$.
  The iterations given by the gradient descent method is
  \[
    x_{n+1} = x_n - \alpha 2 x_n.
  \]
  We need to understand for which values these iterations converge.
\end{example}

\begin{proposition}
  Let $r$ be a strict minimum of $f \in C^2(\R)$.
  Then for all small enough $\alpha > 0$ there exists $\delta > 0$
  such that the gradient descent iterations converge linearly to the minimum
  $r$ for any initial $x_0 \in (r - \delta, r + \delta)$.
\end{proposition}
\begin{proof}
\end{proof}

\subsection{Newton-like method without derivatives*}

\subsection{Polynomials the their roots}

% Fundamental theorem of algebra

% Horner's method

\subsection{Bairstow's method*}

\subsection{Homotopy method*}

% Solving non-linear equations is NP-hard

\newpage

\section{Numerical linear algebra}
\subsection{Solving linear equations}
In this subsection, we will focus on solving a system of linear equations
such as
\[
  Ax = b
\]
where $A \in \R^{n \times n}$ is non-singular.
\begin{itemize}
  \item When $A$ is diagonal we can find a solution in $\O(n)$ operations.
  \item When $A$ is upper triangular or lower triangular we can find a
    solution in $\O(n^2)$ operations.
\end{itemize}
Moreover, if $A = LU$ where $L$ is a lower triangular matrix, and $U$ is
an upper triangular matrix we can solve $Ax = b$ by first solving
$Ux = z$ then solving $Lz = b$.
Since both of the solutions can be found at complexity $\O(n^2)$,
the original system $Ax = b$ can be solved in $\O(n^2)$ operations as well.

\begin{definition}[Permutation matrix]
  A permutation $\sigma \colon \set{1,\dots,n} \to \set{1,\dots,n}$
  induces a linear map $\hat \sigma \colon \R^n \to \R^n$ by
  \[
    \left[\hat \sigma(v) \right]_{\sigma(i)} = v_i
  \]
  for all $1 \le i \le n$.
  The matrix represeting the linear operator $\hat \sigma$ is the matrix
  $P_{\sigma}$ defined by
  \[
    P_{i,j} :=
    \begin{cases}
      1, &i = \sigma(j) \\
      0, &\text{otherwise}
    \end{cases}.
  \]
\end{definition}
\begin{remark}
  When a permutation matrix $P_{\sigma}$ multiplies the matrix $A$ from
  the left, the product $P_{\sigma} A$ is given by
  \[
    \begin{pmatrix}
      \text{---} & R_{\sigma(1)} & \text{---} \\ 
        & \vdots &  \\
      \text{---} & R_{\sigma(n)} & \text{---}
    \end{pmatrix}
  \]
  where $R_i$ is the $i$th row of $A$.
\end{remark}
\begin{remark}
  Every permutation matrix $P$ is orthogonal,
  with its inverse equal to its transpose.
  Moreover, permutation matrices can be characterized as the orthogonal
  matrices whose entries are all non-negative.
\end{remark}

\begin{remark}
  We denote $L^{(r,s,\lambda)}$ the matrix that operates on a matrix $A$
  by adding the $r$th row of $A$ times $\lambda$ to the $s$th row of $A$.
  Formally,
  \[
    L^{(r,s,\lambda)} = I_n + \lambda E^{sl}
  \]
  where $E^{sl}_{sl} = 1$ and otherwise $E_{ij} = 0$.
\end{remark}

\begin{lemma}
  \label{lem:gaussians-commute}
  Let $P_{\sigma}$ be a permutation matrix which swaps the rows $t$ and $u$
  and leaves the rest of the rows fixed.
  Let $L^{r,s,\lambda}$ be with $r > s$,
  and assume that $r < \min\set{u,t}$.
  Then
  \[
    P L^{r,s,\lambda} = L^{r,\sigma(s),\lambda} P
  \]
  and $\sigma(s) > r$.
\end{lemma}
\begin{proof}
  To be added
\end{proof}

\subsubsection{Guassiam elimination}
The process of Guassian elimination on $A$ can be given by first multiplying
$A_0 = A$ by a permutation matrix such that the element in
$(P_0 A_0)_{11} \neq 0$ and then multiplying $P_0 A_0$ on the left
by $n - 1$ lower triangular matrices $L^{1,j,\lambda_j}$ and get a matrix
\[
  A_1 = L^{1,n,\lambda_n} \cdots L^{1,2,\lambda_2} P_0 A_0.
\]
Denote
\[
  L^{k} := L^{k,n,\lambda_n} \cdots L^{1,k+1,\lambda_{k+1}}
\]
and then by the end of the elimination we have
\[
  U = A_{n-1} = L_{n-1} P_{n-2} L^{n-2} P_{n-3} \cdots L^1 P_0 A.
\]
where $U$ is an upper triangular matrix since it is after the elimination.

\begin{theorem}[$LU$ decomposition]
  Let $A \in \R^{n \times n}$ be a non-singular matrix,
  let $P_0,\dots,P_{n-2}$ be the permutation matrices from above,
  let $L^{n-1},\dots,L^1$ be the lower triangular matrices from above,
  and let $U$ be the upper triangular matrix obtained from applying
  Gaussian elimination on $A$.
  Then
  \[
    PA = LU
  \]
  where $L$ is some triangular matrix, and $P$ is the permutation matrix
  $P = P_{n-2} \cdots P_{0}$.
\end{theorem}
\begin{proof}
  Each $P_k$ is a permutation matrix $P_{\sigma_k}$ that permutates
  two rows $u$ and $t$, and each $L^r$ for $r \le k$ is a product of
  matrices of the form $L^{r,s,\lambda}$.
  Accodring to \Cref{lem:gaussians-commute} these types of matrices
  commute (with a slight modification) and thus
  \begin{align*}
    P_k L^r &=
    P_k L^{r,n,\lambda_n} \cdots L^{1,r+1,\lambda_{r+1}} \\ &=
    L^{r,\sigma_k(n),\lambda_n} \cdots L^{1,\sigma_k(r+1),\lambda_{r+1}}
    P_k \\ &=
    L^{r,n,\lambda_n} \cdots L^{1,\sigma_k(r+1),\lambda_{r+1}} P_k
  \end{align*}
\end{proof}

\begin{remark}
  The complexity to find the $LU$ decomposition of a matrix is $\O(n^3)$,
  but this gives us $A = P^T LU$, then we can solve $Ax = P^T LU x = b$
  in complexity $\O(n^2)$.

  This might seem useless at first, because solving $Ax = b$ is still
  $\O(n^3 + n^2) = \O(n^3)$.
  However, if $A$ is fixed and we want to find the solutions $A x = b_i$
  for $b_1,\dots,b_m$, the complexity of solving the equations without
  $LU$ decomposition will be $\O(m n^3)$,
  and with $LU$ decomposition $\O(n^3 + m n^2)$ which is faster in practice.
\end{remark}
\begin{example}
  We can use $LU$ decomposition to find the inverse of $A$ by solving
  the equations
  \[
    A c_j = e_j \quad j = 1,\dots,n
  \]
  in $\O(n^3 + m n^2)$ and then
  \[
    A^{-1} =
    \begin{pmatrix}
      \vert & \vert & \vert \\
      c_1 & \cdots & c_n \\
      \vert & \vert & \vert
    \end{pmatrix}.
  \]
\end{example}
\begin{example}
  We can use $LU$ decomposition to compute the determinant of $A$ with
  \[
    \det(A) = (\det (P))^{-1} \det(L) \det(U).
  \]
  This is very efficient because we can calculate the determinants of
  permutations matrices and triangular matrices in linear time.
\end{example}

\subsection{Scaled row pivoting}
Consider the equation
\[
  \begin{pmatrix}
    0 & 1 \\
    1 & 1
  \end{pmatrix}
  \begin{pmatrix}
    x \\ y
  \end{pmatrix} =
  \begin{pmatrix}
    1 \\ 2
  \end{pmatrix}.
\]
The solution is $(x,y) = (1,1)$.
However, if we modify the equation slightly with a very small $\epsilon > 0$
\[
  \begin{pmatrix}
    \epsilon & 1 \\
    1 & 1
  \end{pmatrix}
  \begin{pmatrix}
    x \\ y
  \end{pmatrix} =
  \begin{pmatrix}
    1 \\ 2
  \end{pmatrix}
\]
After applying Gaussian elimination we get
\[
  \begin{pmatrix}
    \epsilon & 1 \\
    1 & 1 - \epsilon^{-1}
  \end{pmatrix}
  \begin{pmatrix}
    x \\ y
  \end{pmatrix} =
  \begin{pmatrix}
    1 \\ 2 - \epsilon^{-1}
  \end{pmatrix}
\]
and the solutions will be
\begin{align*}
  x &= \epsilon^{-1}\del{1 - \frac{2 - \epsilon^{-1}}{1 - \epsilon^{1}}} =
  - \frac{\epsilon^{-1}}{1 - \epsilon^{-1}} \approx 1 \\
  y &= \frac{2 - \epsilon^{-1}}{1 - \epsilon^{-1}} \approx 1.
\end{align*}
However, when $\epsilon$ is small enough, we have that
\[
  \fl(1 - \epsilon^{-1}) = \fl(2 - \epsilon^{1}) \implies x = 0
\]
which then gives that the solution is $(x,y) = (0,1)$ which is far from
the correct solution.
This shows us that when performing Gaussian elimination we need to have
scaled row pivoting.
With scaled row pivoting, for each row we set
\[
  s_i := \max\set{\abs{a_{i2},\dots,\abs{a_{in}}}}
\]
and then choose the first rwo to be the one for which $\abs{a_{i1}} / s_i$
is maximal in all the steps.

Using this approach in our previous example gives
\[
  \begin{pmatrix}
    1 & 1 \\
    \epsilon & 1
  \end{pmatrix}
  \begin{pmatrix}
    y \\ x
  \end{pmatrix} =
  \begin{pmatrix}
    2 \\ 1
  \end{pmatrix} \implies
  \begin{pmatrix}
    1 & 1 \\
    0 & 1 - \epsilon
  \end{pmatrix}
  \begin{pmatrix}
    y \\ x
  \end{pmatrix} =
  \begin{pmatrix}
    2 \\ 1 - 2 \epsilon
  \end{pmatrix}
\]
and then even for a very small $\epsilon > 0$ we have that
$\fl(1 - \epsilon) = 1$ and $\fl(1 - 2 \epsilon) = 1$,
which leads to $(x,y) = (1,1)$ which is very close to the desired solution.

\subsection{Condition number}
\begin{question}
  We know that we can solve $Ax = b$ when $\det(A) \neq 0$ but what happens
  when $\det(A)$ is very small?
\end{question}
\begin{answer}
  The solution will be very unstable.
\end{answer}

\begin{definition}[Condition number]
  The condition number of a non-singular matrix $A$ is defined as
  \[
    \kappa(A) := \norm{A} \cdot \norm{A^{-1}}.
  \]
\end{definition}
\begin{lemma}
  Let $A$ be a non-singular matrix.
  Then $\kappa(A) > 1$.
\end{lemma}
\begin{proof}
  \[
    1 = \norm{A A^{-1}} \le \norm{A} \cdot \norm{A^{-1}} = \kappa(A).
  \]
\end{proof}

\begin{remark}
  Since we are solving equations numerically we only get approximate
  solutions $\hat x \approx x$.
  Denote $\hat b = A \hat x$.
\end{remark}

\begin{theorem}
  Let $A \in \R^{n \times n}$ be non-singular,
  and let $x, \hat x, b, \hat b$ be nonzero vectors in $\R^n$
  such that
  \[
    Ax = b \tand
    A \hat x = \hat b
  \]
  then
  \[
    \kappa^{-1}(A) \frac{\norm{b - \hat b}}{\norm{b}} \le
    \frac{\norm{x - \hat x}}{\norm{x}} \le
    \kappa(A) \frac{\norm{b - \hat b}}{\norm{b}}.
  \]
\end{theorem}
\begin{proof}
  We have that
  \begin{align*}
    \norm{x - \hat x} &= \norm{A^{-1}(b - \hat b)} \le
    \norm{A^{-1}} \norm{b - \hat b} \\ &=
    \norm{A^{-1}} \norm{b - \hat b} \norm{Ax} \norm{b}^{-1} \le
    \kappa(A) \norm{b - \hat b} \norm{x} \norm{b}^{-1}.
  \end{align*}
  Similarly since $x = A^{-1} b$ and $\hat x = A^{-1} \hat b$ ew get
  \[
    \frac{\norm{b - \hat b}}{\norm{b}} \le
    \kappa(A^{-1}) \frac{\norm{x - \hat x}}{\norm{x}} =
    \kappa(A) \frac{\norm{x - \hat x}}{\norm{x}}
  \]
\end{proof}

\begin{theorem}
  Let $A \in R^{n \times n}$ be a non-singular matrix with an orthonomal
  basis $\mathcal B = \set{v_i}_{i=1}^{n}$ and eigenvalues
  $\Lambda = \set{\lambda_i}_{i=1}^{n}$.
  Then
  \[
    \kappa(A) = \frac{\min \Lambda}{\max \Lambda}.
  \]
\end{theorem}
\begin{proof}
  Since each vector can be decomposed to a linear combination of the basis
  vectors, we get that $\norm{A} = \norm{A}_{\op} = \max \Lambda$.
  Similarly, since the eigenvalues of $A^{-1}$ are $\lambda_i^{-1}$ we
  get that $\norm{A^{-1}} = 1 / \min \Lambda$.
  It follows that
  \[
    \kappa(A) = \frac{\min \Lambda}{\max \Lambda}
  \]
  which completes the proof.
\end{proof}

\subsection{Pure \texorpdfstring{$LU$}{LU} and Cholesky decomposition}

There are some types of matrices such that we do not need any pivoting.

\begin{theorem}
  Let $A \in \R^{n \times n}$ be a matrix whose principle minors
  $A_k \in \R^{k \times k}$ are all non-singular.
  Then there exist a lower triangular matrix $L$ and an upper triangular
  matrix $U$ such that $A = LU$ and the diagonla entries of $L$ are $1$.
\end{theorem}
\begin{proof}
  To be added
\end{proof}

%%%%%%%%%%%%%%%%%%%%% The above are weird

\subsection{The power method}
\begin{example}
  Consider the matrix
  \[
    A = \begin{pmatrix}
      -2 & 0 \\
      0 & 1
    \end{pmatrix}.
  \]
  It clearly has the eigenvalues $\lambda_1 = -2$ and $\lambda_2 = 1$.
  Given an aribitrary vector $x = (x_1, x_2)$ we get
  \[
    A^k x =
    \begin{pmatrix}
      (-2)^k x_1 \\
      x_2
    \end{pmatrix}
  \]
  which in a sense converges to the eigenvector corresponding to $\lambda_1$
  (the vector $e_1$).

  In order to see this more clearly, we first need to normalize the vector
  by defining recursively
  \[
    x^{(k)} = \frac{A^k x}{\norm{A^k x}}.
  \]
  We see that $x^{(2k)}$ converges to $e_1$,
  and $x^{(2k+1)}$ converges to $-e_1$,
  which are scalar multiplies of each other.
  Either way we have that
  \[
    \norm{A x^{(k) - \lambda_1 x^{(k)}}} \taking{k \to \infty} 0.
  \]
  Let $v_1$ be the eigenvector of $\lambda_1$.
  Since for any vector $w \in \R^2$ such that $\ip{v_1}{w} \neq 0$ we get
  \[
    \frac{\ip{A v_1}{w}}{\ip{v_1}{w}} = \lambda_1
  \]
  we can also approximate it by
  \[
    \frac{\ip{A x^{(k)}}{w}}{\ip{x^{(k)}}{w}} \taking{k \to \infty} \lambda_1
  \]
\end{example}

What we have seen in the previous example is called the power method.
It states that given $A \in \C^{n \times n}$ and $x^{(0)}, w \in \C^{n}$,
if we recursively define
\begin{align*}
  x^{(k)} &= \frac{A x^{(n-1)}}{\norm{A x^{(n-1)}}} \\
  \lambda^{(k)} &= \frac{\ip{A x^{(k)}}{w}}{\ip{x^{(k)}}{w}}
\end{align*}
the sequence $x^{(k)}$ will converge to the ``first'' eigenvector
under some conditions.
\begin{remark}
  This means that $\lambda^{(k)}$ will necessarily converge to the
  corresponding eigenvalue.
\end{remark}

\begin{theorem}
  Let $n \geq 2$ be a natrual number,
  let $A \in \C^{n \times n}$,
  let $w^{(0)}, w \in \C^n$,
  and assume that
  \begin{enumerate}
    \item[(1)] $A$ is diagonalizable.
    \item[(2)] $A$ must have a \textit{dominant} eigenvalue 
      which we will denote $\lambda_1$.
    \item[(3)] When $x^{(0)}$ is expressed as a linear combination
      \[
        x^{(0)} = a_1 u_1 + a_2 u_2 + \cdots + a_n u_n
      \]
      the coefficient $a_1$ is not zero.
    \item[(4)] $\ip{w}{u_1} \neq 0$.
  \end{enumerate}
  Then $x^{(k)}$ is well defined for all $k \geq 1$,
  and $\lambda^{(k)}$ is well defined for all sufficiently large $k$,
  and
  \begin{align*}
    &\lambda^{(k)} \taking{k \to \infty} \lambda_1 \\
    &\norm{A x^{(k) - \lambda_1 x^{(k)}}} \taking{k \to \infty} 0.
  \end{align*}
\end{theorem}
\begin{remark}
  A dominant eigenvalue is an eigenvalue $\lambda_1$ such that for
  all $k \geq 2$ we have $\abs{\lambda_1} > \abs{\lambda_k}$.
  This is also the must restricitve condition,
  the rest of the conditions are usually satisfied.
  If the first two eigenvalues have similar magnitude the convergence
  will be very slow.
\end{remark}
\begin{proof}
\end{proof}

\begin{remark}
  In linear algebra we saw that for two column vectors
  $x,y \in \R^{n \times 1}$ we have that the standard inner product
  is
  \[
    \ip{x}{y} = y^T x,
  \]
  but we also find the \textit{outer product} (or tensor product) useful,
  which is defined as
  \[
    x \otimes y = x y^T \in \R^{n \times n}. 
  \]
  This is a linear transformation, such that
  \[
    (x \otimes y) z = \ip{y}{z} x.
  \]
  In particular if $v_1, v_2$ are orthogonal, and of norm one,
  \[
    (v_1 v_1^T) v_1 = v_1
    \tand
    (v_1 v_1^T) v_2 = 0.
  \]
\end{remark}

\subsection{The power method beyond a singlue eigenvalue}
Suppose we have a symmetric real matrix with eigenvalues such that
\[
  \abs{\lambda_1} > \abs{\lambda_2} > \abs{\lambda_3}
  \geq \cdots \geq \abs{\lambda_n}.
\]
After finding the first eigenvalue and eigenvector $(\lambda_1,u_1)$
using the power method, we define $\tilde A = A - \lambda_1 u_1 u_1^T$.
This gives that
\[
  \del{A - \lambda_1 u_1 u_1^T} u_j =
  \lambda_j u_j - \lambda_1 u_1 u_1^T u_j =
  \begin{cases}
    0, &j = 1 \\
    \lambda_j u_j, &j \neq 1
  \end{cases}.
\]
This way the dominant eigenvalue of $\tilde A$ is $\lambda_2$,
and in this way we can theoretically find all eigenvalues of $A$.
\begin{remark}
  To find all eigenvalues of $A$ we need to slightly modify the conditions
  so that
  \[
    \abs{\lambda_1} > \abs{\lambda_2} > \abs{\lambda_3}
    > \cdots > \abs{\lambda_n}.
  \]
\end{remark}
\begin{remark}
  In practice, the power method is used when we need to find only a few
  eigenvectors, in order to find all of them other method, such as
  the $QR$ decomposition method are preferred.
\end{remark}

\begin{remark}
  In order to find the smallest eigenvalue,
  we need to assume that there exists a dominant eigenvalue,
  and a unique minimal eigenvalue, which are all positive.
  Then after finding $\lambda_1$,
  we can define $\tilde A = \lambda_1 I_n - A$,
  which has the eigenvalues $\set{\lambda_1 - \lambda_i}_{i=1}^{n}$.
  The dominant eigenvalue is $\lambda_1 - \lambda_n$, so we can find
  it by using the power method on $\tilde A$.
  Since we already know $\lambda_1$,
  this means we can find $\lambda_n$.
\end{remark}

\subsection{Laplacian matrix}
\begin{definition}[Adjacency matrix]
  Let $G = (V,E)$ be a simple graph.
  We define its adjacency matrix $A$ by
  \[
    (A)_{ij} =
    \begin{cases}
      1, &\set{i,j} \in E \\
      0, &\set{i,j} \notin E
    \end{cases}.
  \]
\end{definition}
\begin{remark}
  The degree $d_i$ of the $i$th vertex is given by
  \[
    d_i = \sum_{j=1}^{n} A_{ij}.
  \]
\end{remark}
\begin{definition}[Laplacian matrix]
  Let $G = (V,E)$ be a graph,
  let $A$ be its adjacency matrix,
  and let $\set{d_i}_{i=1}^{n}$ be the degrees of the vertices of $G$.
  The Laplacian matrix of $G$ is
  \[
    L = \diag(d_1,\dots,d_n) - A.
  \]
\end{definition}







\end{document}
